%% 第 4 章：算子多项式
%% 本文件由 tools/md2tex.py 自动生成，请勿直接编辑。

\chapter{算子多项式}


\begin{jdefn}{算子多项式的定义}{r:5:1}
对于多项式：

\[
p(z)=a_0+a_1z+\cdots+a_mz^m,
\]

定义：

\[
p(T)
=
a_0I+a_1T+\cdots+a_mT^m.
\]

这个定义有意义，因为 $I,T,T^2,\ldots$ 都是从 $V$ 到 $V$ 的线性算子，可以相加与数乘。
\end{jdefn}

\begin{jprop}{算子多项式的乘法}{r:5:2}
若 $p,q$ 是两个多项式，则：

\[
p(T)q(T)=(pq)(T).
\]
\end{jprop}

\begin{jidea}{理解}
所有出现在 $p(T)$ 与 $q(T)$ 中的项，都是同一个算子 $T$ 的幂。因此：

\[
T^rT^s=T^{r+s},
\]

且它们彼此可交换。

这使普通多项式的乘法规则可以直接迁移到算子多项式。
\end{jidea}

\begin{jprop}{多项式算子的核和值域是不变子空间}{r:5:3}
对于任意多项式 $p$：

\[
\ker p(T)
\]

与：

\[
\operatorname{range}p(T)
\]

都是 $T$ 的不变子空间。
\end{jprop}

\begin{jproof}{证明思路}
$T$ 与 $p(T)$ 可交换。

若 $v\in\ker p(T)$，则：

\[
p(T)T(v)=T p(T)(v)=0.
\]

所以 $T(v)$ 仍在 $\ker p(T)$ 中。

若 $v\in\operatorname{range}p(T)$，则存在 $u$ 使：

\[
v=p(T)u.
\]

于是：

\[
T(v)=Tp(T)u=p(T)T(u).
\]

所以 $T(v)$ 仍在 $\operatorname{range}p(T)$ 中。
\end{jproof}
